\section{Monte Carlo Tree Search} \label{sec:mcts}

Monte Carlo Tree Search (MCTS) to rodzina algorytmów przeszukiwania drzewa gry, która łączy systematyczną eksplorację przestrzeni stanów z losowym próbkowaniem wyników rozgrywki. Obszerne omówienie metody przedstawili Browne i in.~\cite{Browne2012}, a nowsze modyfikacje i zastosowania podsumowali Świechowski i in.~\cite{Swiechowski2022}.

\subsection{Ogólny schemat MCTS}

Algorytm buduje drzewo stanów gry przez iteracyjne powtarzanie czterech faz~\cite{Browne2012}:

\begin{enumerate}
\item \textbf{Selekcja} --- poczynając od korzenia, czyli aktualnego stanu gry, algorytm schodzi w dół drzewa według reguły wyboru węzła, która równoważy eksplorację słabo zbadanych gałęzi z eksploatacją dotychczas obiecujących węzłów. Najpopularniejszą taką regułą jest UCB1, na której oparty jest wariant UCT:
\[
\mathrm{UCB1}(v) = \frac{w_v}{n_v} + C\sqrt{\frac{\ln n_{\mathrm{parent}}}{n_v}}
\]
gdzie $w_v$ to suma wygranych w węźle $v$, $n_v$ to liczba wizyt węzła $v$, $n_{\mathrm{parent}}$ to liczba wizyt węzła rodzica, a $C$ to stała eksploracji~\cite{Browne2012}.
\item \textbf{Ekspansja} --- do drzewa dodawany jest jeden lub więcej nowych węzłów potomnych.
\item \textbf{Symulacja} --- od nowo dodanego węzła prowadzona jest losowa symulacja rozgrywki (ang.~\textit{rollout}) aż do końca gry lub limitu głębokości.
\item \textbf{Propagacja wsteczna} --- wynik symulacji jest propagowany wstecz do korzenia, aktualizując statystyki odwiedzonych węzłów.
\end{enumerate}

Po wyczerpaniu budżetu obliczeniowego algorytm zwraca akcję prowadzącą do najczęściej odwiedzanego lub najwyżej ocenianego dziecka korzenia.

\subsection{MCTS w środowisku Hearthstone}

Zastosowanie MCTS w Hearthstone wiąże się z kilkoma wyzwaniami. Santos i in.~\cite{Santos2017} zbadali możliwości standardowego MCTS w tym środowisku, wskazując na ogromną przestrzeń akcji w pojedynczej turze jako główną przeszkodę dla efektywnego budowania drzewa. Częściowa obserwowalność (ukryte karty przeciwnika) i losowe efekty kart uniemożliwiają ponadto budowanie deterministycznego drzewa stanów, co wymaga zastosowania procedury determinizacji lub uśredniania po możliwych stanach ukrytych.


\subsection{Agent PredatorMCTS i hybrydowe podejście Świechowskiego}

Kluczową pracą opisującą zastosowanie MCTS w Hearthstone jest praca Świechowskiego, Tajmajera i Janusza~\cite{swiechowski2018}, w której autorzy zbadali wpływ modeli nadzorowanego uczenia maszynowego na skuteczność agentów MCTS. Punktem wyjścia była własna implementacja MCTS przystosowana do środowisk z częściowo ukrytą informacją i losowymi efektami. Autorzy wykazali, że nawet proste sieci neuronowe, wytrenowane na danych z rozgrywek ludzkich graczy, mogą pełnić rolę funkcji oceny stanów planszy (ang.~\textit{value network}) zamiast symulacji \textit{rolloutów}. Uzupełnienie algorytmu o \textbf{sieć przewodnika} (ang.~\textit{policy network}) kierującą wyborem akcji --- tak, aby MCTS skupiał się na ruchach obserwowanych u doświadczonych graczy --- pozwoliło podnieść skuteczność agenta przy tym samym budżecie obliczeniowym~\cite{swiechowski2018}. Dane treningowe do modeli predykcyjnych zostały zebrane w ramach wcześniejszego wyzwania AAIA'17 Data Mining Challenge~\cite{Janusz2017}.

Na podobnej zasadzie działał agent \textbf{PredatorMCTS} autorstwa Fricka i Akkaya~\cite{dockhorn2018}, który wygrał edycję 2018 ścieżki \textit{premade deck playing track}~\cite{hearthstone_competition}. Agent łączył MCTS z modelem predykcji ruchów przeciwnika, a na podstawie obserwacji aktualnego stanu planszy model szacował, które karty i akcje oponent najprawdopodobniej wykona w kolejnej turze. Włączenie tego modelu do procesu symulacji poprawiło jakość decyzji w momentach, gdy wynik zależy od przewidywanej odpowiedzi przeciwnika.

\subsection{MCTS w kolejnych edycjach konkursu}

MCTS utrzymał silną pozycję w kolejnych latach. W edycji 2019 drugie miejsce w \textit{PDPT} zajął bot García-Sáncheza, Marcosa i Fernándeza-Leiva oparty na MCTS, a czwarte --- bot Le i Bornemanna~\cite{hearthstone_competition}. Choe i Kim~\cite{choe2019} pokazali, że redukcja przestrzeni akcji przez grupowanie sekwencji ruchów w jedną \textit{fazę zagrywania} istotnie zmniejsza współczynnik rozgałęzień drzewa i pozwala MCTS działać efektywniej w turach o dużej liczbie możliwych akcji. W edycji 2020 trzecie miejsce zajął bot Romero i Mory oparty na klasycznym MCTS~\cite{hearthstone_competition}, a MCTS po raz pierwszy nie wygrał wtedy ścieżki \textit{premade deck playing track}, zdominowanej przez lookahead. Mimo tego wcześniejsze zwycięstwa i stabilna obecność w czołówce potwierdzają skuteczność tej klasy metod w Hearthstone.

